%%
%%  INCLUDE: lnre.tex
%%  
%%  some special terminology and notation for lexical statistics and LNRE distributions
%%


%% \onetoS
%% {1,...,S}, the set of population type indices
\newcommand{\onetoS}{\set{1, \dots, S}}

%% \sumtoS, \sumtoS[k]
%% sum over all population types, with index variable i (default) or k
\newcommand{\sumtoS}[1][i]{\sum_{#1 = 1}^S}

%% \sumi, \sumj, \sumk, \sumi[n], ... 
%% the usual summation over all types, with different variables
\newcommand{\sumi}[1][S]{\sum_{i=1}^{#1}} 
\newcommand{\sumj}[1][S]{\sum_{j=1}^{#1}} 
\newcommand{\sumk}[1][S]{\sum_{k=1}^{#1}}

%% \Vmr = V_{m,\rho}, \Vmr[\rho_1], \VmR = \Vmr[> \rho], \Rmr
\newcommand{\Vmr}[1][\rho]{V_{m,#1}}
\newcommand{\VmR}[1][\rho]{V_{m,> #1}}
\newcommand{\Rmr}[1][\rho]{R_{m,#1}}

%% union vocabulary and frequency spectrum
\newcommand{\Vcup}[1][A\cup B]{V^{#1}}
\newcommand{\Vcap}[1][A\cap B]{V^{#1}}
\newcommand{\VdiffA}[1][A\setminus B]{V^{#1}}
\newcommand{\VdiffB}[1][B\setminus A]{V^{#1}}



%%% Local Variables: 
%%% mode: plain-tex
%%% TeX-master: t
%%% End: 
